% from ega3-1-fr.tex
\oldpage[0\textsubscript{III}]{59}
qu'on pourrait identifier aux «~edge-homomorphismes~» de la suite spectrale
de Leray (cf. (\hyperref[0.12.1.7-fr]{12.1.7})). Il suffit d'opérer sur les
préfaisceaux auxquels sont associés les faisceaux images supérieures
(\hyperref[0.12.2.1-fr]{12.2.1}). Pour cela, considérons un ouvert quelconque
$W$ de $Z$ et son image réciproque $g^{-1}(W)$ dans $Y$; on a un
di-homomorphisme canonique
\[
\label{0.12.2.5.3-fr}
\tag{12.2.5.3}
H^n(g^{-1}(W),f_*(\mathcal F))\longrightarrow
H^n(f^{-1}(g^{-1}(W)),f^*(f_*(\mathcal F)))
\]
les anneaux correspondants étant $\Gamma(g^{-1}(W),\mathcal O_Y)$ et
$\Gamma(h^{-1}(W),\mathcal O_X)$; d'autre part, l'homomorphisme canonique
(\hyperref[0.4.4.3.3-fr]{0\textsubscript{I}, 4.4.3.3}) fournit par
fonctorialité des homomorphismes canoniques
\[
\label{0.12.2.5.4-fr}
\tag{12.2.5.4}
H^n(h^{-1}(W),f^*(f_*(\mathcal F)))\longrightarrow
H^n(h^{-1}(W),\mathcal F)
\]
qui sont des homomorphismes de $\Gamma(h^{-1}(W),\mathcal O_X)$-modules.
Compte tenu de l'homomorphisme d'anneaux
$\Gamma(W,\mathcal O_Z)\to\Gamma(g^{-1}(W),\mathcal O_Y)$, on voit qu'en
composant (\hyperref[0.12.2.5.4-fr]{12.2.5.4}) et
(\hyperref[0.12.2.5.3-fr]{12.2.5.3}), on obtient un homomorphisme de
préfaisceaux, qui fournit l'homomorphisme de faisceaux
(\hyperref[0.12.2.5.1-fr]{12.2.5.1}).

La définition de (\hyperref[0.12.2.5.2-fr]{12.2.5.2}) est encore plus
simple; par définition, $R^n h_*(\mathcal F)$ est associé au préfaisceau
$W\leadsto H^n(f^{-1}(g^{-1}(W)),\mathcal F)$ et $R^n f_*(\mathcal F)$ au
préfaisceau $V\leadsto H^n(f^{-1}(V),\mathcal F)$; on a donc un homomorphisme
canonique
\[
H^n(f^{-1}(g^{-1}(W)),\mathcal F)\longrightarrow
\Gamma(g^{-1}(W),R^n f_*(\mathcal F)),
\]
et il est immédiat que ces homomorphismes définissent un homomorphisme de
préfaisceaux, qui à son tour définit
(\hyperref[0.12.2.5.2-fr]{12.2.5.2}).
\end{env}

\begin{env}[12.2.6]
\label{0.12.2.6-fr}
Sous les hypothèses de (\hyperref[0.12.2.4-fr]{12.2.4}), soient
$\mathcal F$, $\mathcal G$, $\mathcal K$ trois $\mathcal O_X$-Modules et
$u:\mathcal F\otimes\mathcal G\to\mathcal K$ un
$\mathcal O_X$-homomorphisme. On a alors des diagrammes commutatifs
\[
\xymatrix@C=10pt@R=24pt{
R^p g_*(f_*(\mathcal F))\otimes_{\mathcal O_Z}R^q g_*(f_*(\mathcal G))
\ar[r]\ar[d]&R^{p+q}g_*(f_*(\mathcal K))\ar[d]\\
R^p h_*(\mathcal F)\otimes_{\mathcal O_Z}R^q h_*(\mathcal G)
\ar[r]&R^{p+q}h_*(\mathcal K)
}
\tag{12.2.6.1}\label{0.12.2.6.1-fr}
\]
et
\[
\xymatrix@C=10pt@R=24pt{
R^p h_*(\mathcal F)\otimes_{\mathcal O_Z}R^q h_*(\mathcal G)
\ar[r]\ar[d]&R^{p+q}h_*(\mathcal K)\ar[d]\\
g_*(R^p f_*(\mathcal F))\otimes_{\mathcal O_Z}g_*(R^q f_*(\mathcal G))
\ar[r]&g_*(R^{p+q}f_*(\mathcal K))
}
\tag{12.2.6.2}\label{0.12.2.6.2-fr}
\]
